% !TeX root = ../thuthesis-example.tex

% 中英文摘要和关键字
\begin{abstract}
本文围绕国际中文学习平台 LingoBridge 的教学应用场景，设计并实现了一套基于 DKT 知识追踪与 Levenshtein 距离校准的国际中文自适应学习推荐系统。系统以教师端、学生端和管理端三类业务入口为基础，采集课程、作业、录音练习、批改反馈、学习时长、完成率和班级排名等多维学习数据，经过数据清洗、特征构建和指标聚合，形成可用于学习状态诊断与练习推荐的行为画像。近年研究表明，学习分析能够通过学生行为数据、反馈循环、能力分层和可视化洞察支持个性化学习，人工智能也正在高等教育个性化学习场景中发挥推荐、监测和自适应反馈作用\cite{Ref-LearningAnalyticsPersonalized2024,Ref-AIPersonalizedLearningHE2025}。与单纯展示课程内容的传统教学平台不同，本系统将人工智能模块嵌入真实应用流程，使学生能够获得针对拼音、声调、听读和作业完成情况的个性化练习建议，教师能够识别学习风险和班级薄弱环节，管理端能够作为数据中台监控平台整体运行与教学效果。

在算法设计上，系统采用“DKT 知识追踪 + Levenshtein 距离校准 + 可解释推荐”的轻量化方案。首先对学习记录进行去重、缺失值处理和时间窗口聚合，构建练习频率、提交稳定性、错题分布和学习进步趋势等特征；其次通过 DKT 序列模型估计学生对拼音、声调和词汇知识点的动态掌握状态；最后利用 Levenshtein 编辑距离对 ASR 字幕或拼音转写偏误进行相似度校准，并结合薄弱知识点生成练习推荐。教育推荐系统研究也指出，推荐效果不应只看相关性，还应关注学习结果、评价方法和教学影响\cite{Ref-EducationalRecommenderEvaluation2024}。该方案兼顾课程大作业对人工智能方法、数据分析能力和工程实现能力的要求，也避免了在本地演示环境中过度依赖不可解释的大模型。实验与演示结果表明，新增 AI 模块能够在不破坏 LingoBridge MVP 主流程的前提下，增强三端首页、个人排名页和学生总览页的数据分析表达，使平台从“教学管理工具”进一步演进为具备智能诊断与推荐能力的学习辅助系统。  \thusetup{
    keywords = {DKT知识追踪,Levenshtein距离,自适应推荐,国际中文教育,学习行为分析},
  }
\end{abstract}

\begin{abstract*}
This paper presents the design and implementation of an adaptive international Chinese learning recommendation system based on Deep Knowledge Tracing (DKT) and Levenshtein distance calibration for LingoBridge, an international Chinese learning platform. Built on the existing teacher, student, and administrator workflows, the system collects multi-dimensional learning data, including courses, assignments, audio practice records, teacher feedback, learning duration, completion rates, and class rankings. After data cleaning, feature construction, and metric aggregation, these records are transformed into learning profiles that support diagnosis and recommendation. Recent studies show that learning analytics can support personalized learning through behavioral data, feedback loops, learner classification, and classroom-level visualization, while artificial intelligence is increasingly used in higher education for recommendation, monitoring, and adaptive feedback\cite{Ref-LearningAnalyticsPersonalized2024,Ref-AIPersonalizedLearningHE2025}. Unlike a conventional teaching platform that only displays learning materials, the proposed system embeds artificial intelligence modules into the actual application flow: students receive personalized practice suggestions for pinyin, tones, listening, reading, and homework; teachers identify at-risk learners and class-wide weaknesses; and administrators monitor global learning and platform indicators through a data-center-style dashboard.

From the algorithmic perspective, the system follows a lightweight and interpretable strategy that combines DKT, Levenshtein distance calibration, and explainable recommendation. Learning records are first deduplicated, normalized, and aggregated over time windows. DKT is then used to estimate students' dynamic mastery of pinyin, tone, and vocabulary knowledge points, while Levenshtein distance is used to calibrate ASR subtitles or pinyin transcription errors through string similarity matching. Based on the resulting mastery state and corrected text evidence, the system recommends exercises according to weak knowledge points, historical performance, and class-level comparisons. Research on educational recommender systems also suggests that recommendation quality should be evaluated not only by relevance, but also by learning outcomes, assessment methods, and pedagogical impact\cite{Ref-EducationalRecommenderEvaluation2024}. This design satisfies the course project requirements for artificial intelligence methods, data analysis, interaction design, and engineering implementation, while avoiding excessive dependence on opaque large models in a local demonstration environment. The experimental demonstration shows that the AI modules strengthen the data-analysis and recommendation capabilities of the student dashboard, personal ranking page, teacher overview, and administrator console without disrupting the original LingoBridge MVP workflow.
  \thusetup{
    keywords* = {Deep Knowledge Tracing, Levenshtein Distance, Adaptive Recommendation, International Chinese Education, Learning Analytics},
  }
\end{abstract*}
